\SourceLeaf{060}
\begin{center}
\pdfbookmark[2]{42.}{article-042}42.
\end{center}
Tabulam, qualem hactenus descripsimus, extensione idonea construximus, operique huic adiecimus (Tab. I.). Ad ellipsin sola pars prior spectat; partem alteram, quae motum hyperbolicum complectitur, infra explicabimus. Argumentum tabulae, quod est quantitas $A$, per singulas partes millesimas a $0$ usque ad $0{,}300$ progreditur; sequuntur $\log B$ et $C$, quas quantitates in partibus $10\,000\,000^{\mathrm{mis}}$, sive ad septem figuras decimales expressas subintelligere oportet; cifrae enim primae, figuris significativis praeeuntes, suppressae sunt; columna denique quarta exhibet quantitatem $T$ primo ad 5 dein ad 6 figuras computatam, quae praecisio abunde sufficit, quum haec columna ad eum tantummodo usum requiratur, ut argumento $T$ valores respondentes ipsius $\log B$ et $C$ habeantur, quoties ad normam art. praec. $t$ ex $v$ determinare lubet. Quum problema inversum, quod longe frequentioris usus est, puta determinatio ipsius $v$ et $r$ ex $t$, omnino absque quantitatis $T$ subsidio absolvatur, quantitatem $A$ pro argumento tabulae nostrae eligere maluimus quam $T$, quae alioquin argumentum aeque fere idoneum fuisset, imo tabulae constructionem aliquantulum facilitavisset. Haud superfluum erit monere, omnes tabulae numeros ad decem figuras ab origine calculatos fuisse, septemque adeo figuris, quas hic damus, ubique tuto confidere licere; methodis autem analyticis ad hunc laborem in usum vocatis hoc loco immorari non possumus, quarum explicatione copiosa nimium ab instituto nostro distraheremur. Ceterum tabulae extensio omnibus casibus, ubi methodum hactenus expositam sequi prodest, abunde sufficit, quum ultra limitem $A=0{,}3$, cui respondet $T=0{,}392374$ sive $E=64^0\,8'$, methodis artificialibus commode ut supra ostensum est abstinere liceat.
